\chapter{Procedure operative}

\section{Avvio con immagini GHCR}
Dalla directory \texttt{dashboard}:

\begin{lstlisting}
docker compose -f docker-compose.prod.yml pull
docker compose -f docker-compose.prod.yml up -d
docker compose -f docker-compose.prod.yml ps
\end{lstlisting}

Avviare poi il simulatore C++ dalla root del progetto:

\begin{lstlisting}
./hardware/plantpulse
\end{lstlisting}

La dashboard è raggiungibile su \url{http://localhost:8080}; l'API è su \url{http://127.0.0.1:8000} e la documentazione FastAPI su \url{http://127.0.0.1:8000/docs}.

\section{Arresto e riavvio}
\begin{lstlisting}
docker compose -f docker-compose.prod.yml stop
docker compose -f docker-compose.prod.yml start
\end{lstlisting}

Per ricreare i container mantenendo i dati, è possibile usare \texttt{down} senza \texttt{-v}, seguito da \texttt{up -d}.

\section{Aggiornamento delle immagini}
Dopo una nuova pubblicazione da GitHub Actions:

\begin{lstlisting}
docker compose -f docker-compose.prod.yml pull
docker compose -f docker-compose.prod.yml up -d
\end{lstlisting}

\section{Controlli rapidi}
\begin{lstlisting}
docker compose -f docker-compose.prod.yml ps
docker compose -f docker-compose.prod.yml logs -f backend
curl http://127.0.0.1:8000/
curl http://127.0.0.1:8000/measurements/latest
\end{lstlisting}

Per i test backend, dalla directory \texttt{dashboard/backend}, attivare l'ambiente virtuale e lanciare \texttt{pytest -v}. Per il frontend, dalla directory \texttt{dashboard/frontend}, eseguire \texttt{npm ci}, \texttt{npm run lint} e \texttt{npm run build}.

\section{Troubleshooting}
\begin{itemize}
    \item Database non pronto: controllare lo stato e i log del container PostgreSQL.
    \item Errore di connessione backend/database: controllare il file \texttt{dashboard/.env}, le credenziali e il valore \texttt{DB\_HOST=db} nel container backend.
    \item Dashboard senza dati: verificare che il simulatore sia in esecuzione e che le richieste POST siano accettate.
    \item Comandi in attesa: verificare che il simulatore stia interrogando l'endpoint dei comandi e consultare la cronologia.
\end{itemize}
